% Gemeinsame Originalaussagen anzeigen, ohne sie erneut zu registrieren.
\newcommand{\ReconstructionProofStatement}[1]{%
  \par\Needspace{26\baselineskip}%
  \hypertarget{reconstruction.proof.#1}{}%
  \begingroup
  \RenewDocumentCommand{\FormulaThmDeltaK}{o m m +m}{%
    \noindent\textbf{Beweis zu \FormulaRefAuto{##3}[thm]}%
    \IfNoValueF{##1}{\space(##1)}\par
    {\small\begin{DeltaContext}{ }##4\end{DeltaContext}}%
    {\ifcsname ReconstructionCompactStatement@##3\endcsname\small\fi
      \[##2\]}%
  }%
  \ReconstructionStatement{#1}%
  \endgroup
}

% Diese lange Originalaussage braucht etwas kleineren Formelsatz.
\expandafter\newcommand\csname ReconstructionCompactStatement@LargePowerUnitSubsetImage\endcsname{}

% Die einzige zu breite Tabellenzeile wird an ihrer vorhandenen Konjunktion
% umbrochen. Ihr Originalkörper und sämtliche Ableitungsdaten bleiben gleich.
\long\def\ReconstructionAlignedConjunction#1\land#2\ReconstructionConjunctionEnd{%
  \begin{aligned}[t]& #1\\& {}\land #2\end{aligned}%
}
\newcommand{\ReconstructionConjoinedStatementRow}[3]{%
  \ReconstructionOriginalStatementRow{#1}{%
    \ReconstructionAlignedConjunction#2\ReconstructionConjunctionEnd}{#3}%
}
\newcommand{\ReconstructionUnitSetImageLayout}{%
  \let\ReconstructionOriginalStatementRow\mxProofStatementRow
  \renewcommand{\mxProofStatementRow}[3]{%
    \ifnum\value{proofstepnr}=20\relax
      \expandafter\ReconstructionConjoinedStatementRow
    \else
      \expandafter\ReconstructionOriginalStatementRow
    \fi{##1}{##2}{##3}%
  }%
}

% Die lange quantifizierte Äquivalenz erhält einen eigenen Quantorkopf.
\long\def\ReconstructionAlignedQuantifier#1\,(#2\ReconstructionQuantifierEnd{%
  \begin{aligned}[t]& #1\\& \quad(#2\end{aligned}%
}
\newcommand{\ReconstructionQuantifiedStatementRow}[3]{%
  \ReconstructionOriginalStatementRow{#1}{%
    \ReconstructionAlignedQuantifier#2\ReconstructionQuantifierEnd}{#3}%
}
\newcommand{\ReconstructionUnitFamilyLayout}{%
  \let\ReconstructionOriginalStatementRow\mxProofStatementRow
  \renewcommand{\mxProofStatementRow}[3]{%
    \ifnum\value{proofstepnr}=19\relax
      \expandafter\ReconstructionQuantifiedStatementRow
    \else
      \expandafter\ReconstructionOriginalStatementRow
    \fi{##1}{##2}{##3}%
  }%
}

% Teilüberschrift und die ersten Zeilen der folgenden Tabelle zusammenhalten.
\let\ReconstructionOriginalProofPartWide\proofpartwide
\renewcommand{\proofpartwide}[1]{%
  \par\Needspace{6\baselineskip}%
  \ReconstructionOriginalProofPartWide{#1}%
}

% Auch der intern benannte Produktzeuge bleibt ausschließlich in Band 40
% registriert. Die unveränderte Tabellenquelle benutzt hier nur die Anzeige.
\RenewDocumentCommand{\mxProofPartOpen}{m o m o o}{%
  \par\medskip
  \noindent
  \IfNoValueTF{#2}{\textbf{Teil}}{\textbf{#2}}%
  \IfNoValueF{#5}{\hfill\FormulaRefAuto{#5}[thm]}\par
  \[#3\]
  \par\small
  \setcounter{proofstepnr}{0}%
  \setlength{\tabcolsep}{2pt}%
  \renewcommand{\arraystretch}{1.15}%
  \begin{longtable}{#1}%
}
